\chapter{Paleomagnetics\\ \large Magnetics Extension}\label{ch:paleomag}

\section{Introduction}

Paleomagnetism is the study of the magnetic record preserved in rocks, sediments, and archaeological materials. It provides fundamental constraints on plate tectonics, geomagnetic field behaviour, and the evolution of Earth's interior. Magnetic signals are recorded as natural remanent magnetization (NRM), which reflects both the geomagnetic field and subsequent geological processes.

The geomagnetic field is approximately dipolar at Earth's surface and can be described by declination ($D$), inclination ($I$), and intensity.

\section{Types of Magnetism}
Magnetic behaviour arises from electron spin and orbital angular momentum.  There are three dominant types of magnetization including diamagnetism, paramagnetism, and ferromagnetism. 
\begin{description}
    \item[diamagnetism,] very weak, negative susceptibility. Induced magnetization opposes the applied field;
    \item[paramagnetism,] weak, positive susceptibility. Magnetization disappears when the field is removed;
    \item[ferromagnetism,] strong magnetization due to parallel alignment of magnetic moments. Capable of permanent magnetization; 
    \item[antiferromagnetism,] antiparallel alignment of unequal moments produces a net magnetization. This is the most important form in paleomagnetism. Adjacent spins align antiparallel, producing no net moment; and
    \item[ferrimagnetism,] antiparallel alignment of unequal moments produces a net magnetization. This is the most important form in paleomagnetism.
\end{description}
All materials are diamagnetic and some are paramagnetic but the effects are small enough that they can be ignored for magnetic surveys.

\section{Rock Magnetism}

Rock magnetism describes how magnetic minerals acquire and retain magnetization.

\subsection{Magnetic Domains}

Magnetic minerals are divided into domains to minimize magnetic energy:
\begin{itemize}
\item \textbf{Single-domain (SD)}: uniform magnetization, highly stable
\item \textbf{Pseudo-single-domain (PSD)}: intermediate behaviour
\item \textbf{Multi-domain (MD)}: subdivided domains, less stable remanence
\end{itemize}

SD grains are the most reliable paleomagnetic recorders.

\section{Magnetic Properties Relevant to Paleomagnetism}

In the previous chapter, we treated magnetization primarily as an induced response to the present-day Earth's magnetic field. This approximation is useful for many magnetic surveys but does not capture the full complexity of how rocks acquire and retain magnetization.

Paleomagnetism focuses on remanent magnetization: the ability of rocks to record and preserve information about past magnetic fields. Understanding this process requires a more detailed treatment of magnetic minerals, domain structure, and the mechanisms of magnetization acquisition.

\subsection{Types of magnetic behaviour}

Magnetic behaviour arises from the alignment of atomic magnetic moments. Materials are classified as:
\begin{itemize}
    \item diamagnetic,
    \item paramagnetic,
    \item ferromagnetic,
    \item antiferromagnetic,
    \item ferrimagnetic.
\end{itemize}

Among these, ferrimagnetic minerals (e.g., magnetite) are the most important for paleomagnetism because they can carry strong and stable remanent magnetization.

Antiferromagnetic minerals (e.g., hematite) may also carry remanence, although typically weaker, while paramagnetic and diamagnetic materials do not retain permanent magnetization.

\subsection{Magnetic domains and grain size}

Magnetic minerals are divided into domains, regions within which magnetic moments are uniformly aligned. Domain structure depends strongly on grain size:

\begin{itemize}
    \item \textbf{Single-domain (SD)} grains: uniformly magnetized; capable of carrying strong and stable remanence.
    \item \textbf{Pseudo-single-domain (PSD)} grains: intermediate behaviour.
    \item \textbf{Multi-domain (MD)} grains: subdivided into many domains; remanence is less stable.
\end{itemize}

The stability of remanent magnetization depends critically on domain state. Fine-grained materials tend to preserve magnetic signals over geological timescales, whereas coarse-grained materials are more easily remagnetized.

This distinction is critical because susceptibility primarily reflects the ease of induced magnetization, while remanence reflects the ability to retain magnetization. As a result,
\begin{center}
\emph{high susceptibility does not necessarily imply strong remanent magnetization, and vice versa.}
\end{center}

\subsection{Remanent magnetization and blocking temperature}

Remanent magnetization is acquired when magnetic minerals become aligned with an external field and subsequently ``lock in'' that alignment.

This process is governed by blocking temperature ($T_b$), which is analogous to but distinct from Curie temperature. While $T_C$ defines the loss of magnetic ordering, $T_b$ defines the temperature below which a particular grain can retain a stable remanent magnetization.

Because blocking temperature depends on grain size, domain state, and time, different grains within a rock may record magnetization over a range of temperatures.

\subsection{Mechanisms of remanence acquisition}

Several mechanisms by which rocks acquire remanent magnetization are important:

\begin{itemize}
    \item \textbf{Thermoremanent magnetization (TRM):} acquired during cooling through the Curie temperature.
    \item \textbf{Detrital remanent magnetization (DRM):} acquired during sediment deposition.
    \item \textbf{Chemical remanent magnetization (CRM):} acquired during mineral growth or alteration.
\end{itemize}

These processes allow rocks to record the direction and, in some cases, the intensity of the Earth's magnetic field at the time of formation.

\subsection{Implications for magnetic interpretation}

Remanent magnetization can significantly affect magnetic anomalies when its magnitude is comparable to or larger than induced magnetization. In such cases:
\begin{itemize}
    \item anomaly shapes may be non-intuitive,
    \item anomaly peaks may be displaced,
    \item reduction-to-the-pole assumptions may fail.
\end{itemize}

Understanding remanence is therefore essential for interpreting complex magnetic data and forms the basis for using magnetism to study plate motions and the history of the Earth's magnetic field.

\subsection{Hysteresis}

Magnetic behaviour is characterized by hysteresis loops:
\begin{itemize}
\item Saturation magnetization ($M_s$)
\item Remanent magnetization ($M_r$)
\item Coercive field ($H_c$)
\end{itemize}

These parameters provide information on grain size and mineralogy.

\subsection{Curie Temperature}

The Curie temperature ($T_C$) is the temperature above which a material loses its permanent magnetization. For magnetite, $T_C \approx 580^\circ$C.

\section{Types of Remanent Magnetization}

Rocks may gain remnant magnetization for a long list of reasons including but not limited to
\begin{description}
    \item[Chemical Remanent Magnetization (CRM),] recrystalization of magnetic minerals or metamorphic alteration which produces magnetic minerals acquire magnetization below the Curie point; 
    \item[Thermoremanent Magnetization (TRM),] magnetization acquired as a rock is cooled across the Curie temperature;
    \item[Detrital Remanent Magnetization (DRM),] grains already magnetized are oriented in the Earth's field as they are deposited or subsequent to deposition;
    \item[Viscous Remanent Magnetization (VRM),] is a secondary magnetization resulting from exposure to week magnetic fields after formation and may be due to heating below the Curie temperature and/or metamorphism; and
    \item[Isothermal Remanent Magnetization (IRM),] acquired quickly at constant temperature (generally due to lightning strikes).
\end{description}

\section{Measurement Techniques}

\subsection{Field Sampling}
Oriented cores are collected using drills. Orientation is determined using magnetic or sun compasses.

\subsection{Laboratory Measurements}
Magnetization is measured using:
\begin{itemize}
\item Spinner magnetometers
\item SQUID magnetometers
\end{itemize}

\subsection{Demagnetization}
Used to isolate primary remanence:
\begin{itemize}
\item Thermal demagnetization
\item Alternating field (AF) demagnetization
\end{itemize}

\subsection{Directional Analysis}
Zijderveld diagrams and Fisher statistics are used to estimate mean directions and uncertainties.

\section{Geomagnetic Field Behaviour}

\subsection{Geocentric Axial Dipole}

The time-averaged geomagnetic field approximates a geocentric axial dipole:
\begin{equation}
\tan I = 2 \tan \lambda
\end{equation}

where $\lambda$ is paleolatitude.

\subsection{Paleosecular Variation}
Short-term variations in field direction and intensity occur on $10^3$--$10^4$ year timescales.

\section{Geomagnetic Reversals}

Geomagnetic reversals involve switching of magnetic polarity.

\subsection{Characteristics}
\begin{itemize}
\item Irregular frequency (typically $10^5$--$10^6$ years)
\item Transition durations of several thousand years
\item Transitional fields often non-dipolar
\end{itemize}

\subsection{Origin}
Reversals are linked to instabilities in the geodynamo within the outer core.

\section{Geomagnetic Polarity Timescale (GPTS)}

The GPTS is a chronology of polarity reversals.

\subsection{Construction}
Based on:
\begin{itemize}
\item Radiometric dating
\item Marine magnetic anomalies
\end{itemize}

Notable features include superchrons, such as the Cretaceous Normal Superchron.

\section{Oceanic Magnetic Anomalies}

\subsection{Magnetic Striping}
Oceanic crust records alternating polarity as it forms at mid-ocean ridges.

\subsection{Vine--Matthews--Morley Hypothesis}
Magnetic anomalies reflect geomagnetic reversals recorded during seafloor spreading.

\subsection{Applications}
\begin{itemize}
\item Determination of spreading rates
\item Dating oceanic lithosphere
\end{itemize}

\section{Plate Motions}

\subsection{Paleolatitude}
Paleolatitude is determined from inclination:
\begin{equation}
\lambda = \arctan\left(\frac{1}{2}\tan I\right)
\end{equation}

\subsection{Apparent Polar Wander}
Changes in paleomagnetic poles through time reflect plate motion.

\subsection{Limitations}
\begin{itemize}
\item No direct longitude constraint
\item Inclination shallowing in sediments
\end{itemize}

\section{Plate Reconstructions}

\subsection{Methods}
Plate motions are reconstructed using:
\begin{itemize}
\item Euler rotations
\item Magnetic anomaly data
\item Geological correlations
\end{itemize}

\subsection{Applications}
\begin{itemize}
\item Supercontinent reconstructions (e.g., Pangaea)
\item Paleoclimate studies
\end{itemize}

\section{Limitations and Challenges}

\begin{itemize}
\item Remagnetization
\item Mineral alteration
\item Non-dipole field effects
\item Dating uncertainties
\end{itemize}

\section{Summary}


